% Generated draft: requires mathematical and source review.
\documentclass[11pt]{article}
\usepackage[margin=25mm]{geometry}
\usepackage{fontspec}
\setmainfont{Latin Modern Roman}
\usepackage{amsmath,amssymb,mathtools,unicode-math}
\setmathfont{Latin Modern Math}
\usepackage{enumitem,needspace,xurl,hyperref}
\setlength{\parindent}{0pt}
\setlength{\parskip}{5pt}
\setlength{\emergencystretch}{3em}
\newcommand{\N}{\mathbb{N}}
\newcommand{\Z}{\mathbb{Z}}
\newcommand{\Q}{\mathbb{Q}}
\newcommand{\R}{\mathbb{R}}
\newcommand{\C}{\mathbb{C}}
\newcommand{\F}{\mathbb{F}}
\newcommand{\A}{\mathbb{A}}
\newcommand{\PP}{\mathbb{P}}
\newcommand{\E}{\mathbb{E}}
\newcommand{\eps}{\varepsilon}
\newcommand{\dd}{\,\mathrm{d}}
\begin{document}
\section*{UnsolvedMath ID: 30000410}
\textbf{Source problem code:} catalogue entry
\subsection{Definitions and mathematical statement}
Let $R=\mathcal O_{F'}$ be the global differential ring in the source, with derivation $\partial_F$, and let $(M,\partial_M)$ be its differential module. Thus $M$ is a finite-rank module in the source's category and
\[
 \partial_M(am)=\partial_F(a)m+a\partial_M(m).
\]
Write $k=\ker(\partial_F)$ for the constant field, interpreted in the associated fraction field when appropriate. Solutions mean horizontal local sections, satisfying $\partial_M s=0$.

The hypothesis is vanishing $p$-curvature at almost all reduction primes, meaning outside a finite exceptional set where the chosen integral model may have bad reduction. In connection notation the characteristic-$p$ curvature is
\[
 \psi_p(D)=\nabla_D^{\,p}-\nabla_{D^p};
\]
it is vanishing, not merely nilpotence, that the source assumes. One must not replace this formula by $\nabla_D^p$ unless the relevant $D^p$ term vanishes.

The report's Conjecture 2 asks whether, near every nonsingular degree-one prime, the local solutions are given by locally bounded power series over $k$. At a degree-one place with residue field identified with $k$ and local parameter $z$, the intended expansions have coefficients in $k$ in a suitable local trivialization. Both the degree-one restriction and the source's local-boundedness requirement are part of the target; no assertion at every higher-degree or singular place is added.

The extract does not define the global differential-ring model, its reduction primes, the precise meaning of degree one, or the arithmetic/valuation criterion called locally bounded. It also does not state its exact local solution-space and trivialization conventions. Ordinary complex-analytic boundedness on a small disk must not be silently substituted for the intended arithmetic boundedness, nor are reduction primes automatically identified with the local places in the conclusion.

The literature assessment checked on 21 August 2026 describes a related Eisenstein algebraicity converse from differential finiteness and global boundedness as a consequence of the unresolved $p$-curvature conjecture, and reports only the algebraic-to-necessary-conditions direction established. It supplies no proof of the local boundedness assertion above. These related status claims are dated source triage, not a new equivalence or an independently verified resolution.
\subsection{Short English statement}
Does almost-everywhere vanishing $p$-curvature force the report's arithmetic boundedness for local power-series solutions at every nonsingular degree-one place? The precise global model and boundedness convention are missing and cannot be replaced by ordinary local analyticity.
\subsection{Sources}
\begin{itemize}
\item \emph{The Arithmetic of Fields} (2006), Oberwolfach Reports, DOI 10.4171/owr/2006/06, pp. 317--378, Conjecture 2 according to the supplied verification, \url{https://doi.org/10.4171/owr/2006/06}.
\item Literature assessment checked 21 August 2026: original report, \url{https://doi.org/10.4171/owr/2006/06}; ``primary report text,'' \url{https://ems.press/content/serial-article-files/46036}.
\item Local UnsolvedMath registry, permanent ID $30000410$, exact source code \texttt{catalogue entry}, created 25 August 2026 and updated 16 September 2026. Supplied verification identifies a single conjecture including the degree-one and local-boundedness qualifications.
\end{itemize}
\end{document}
